% Generic fermionic Gaussian pairing combinatorics before QED specialization.
自由高斯理论的高阶关联函数由二点函数决定。玻色场交换两个算符不会改变符号；费米场每交换两个 Grassmann 奇算符都会产生一个负号。因此一般高阶费米关联函数的每一种完全配对都必须带上把原有算符顺序重排到该配对顺序所需的置换宇称。后面的全同费米子交换号与闭合费米子圈负号都来自这一同一代数结构。


设按原始时间序列排列的 $2n$ 个自由费米型场算符统一记为
\[
\Psi_1,\Psi_2,\ldots,\Psi_{2n},
\]
其中每个 $\Psi_r$ 可以是 $\psi$ 或 $\bar\psi$，并把二点缩并记为
\begin{equation}
 C_{rs}\equiv
 \langle0|\Tord\{\Psi_r\Psi_s\}|0\rangle_{\rm contraction}.
 \label{eq:generic-fermion-contraction-dp8}
\end{equation}
只有场类型与量子数允许的二点缩并才非零。令 $\mathfrak P_{2n}$ 表示集合 $\{1,\ldots,2n\}$ 的全部完全配对；对每个配对 $P$，$(-1)^{\pi(P)}$ 定义为把原始算符序列重排成各配对相邻时所产生的费米置换符号。自由高斯费米场的一般 $2n$ 点函数于是为
\begin{equation}
 \boxed{
 \langle0|\Tord\{\Psi_1\Psi_2\cdots\Psi_{2n}\}|0\rangle
 =\sum_{P\in\mathfrak P_{2n}}
 (-1)^{\pi(P)}
 \prod_{(r,s)\in P} C_{rs}.}
 \label{eq:fermion-wick-general-dp73}
\end{equation}
这一“任意高阶高斯关联函数等于所有二点配对之和，并由费米置换宇称定号”的结果通常称为费米 Wick 定理。其母结构是\eqnrefs{eq:fermion-wick-general-dp73}，低阶公式只是代入后的检验。

例如取 $n=2$ 且算符次序为 $\psi_1\bar\psi_2\psi_3\bar\psi_4$，两个允许的完全收缩具有相反宇称，因此
\begin{equation}
 \boxed{
 \langle0|\Tord\{\psi_1\bar\psi_2\psi_3\bar\psi_4\}|0\rangle
 =S_{12}S_{34}-S_{14}S_{32},}
 \label{eq:fermion-wick-four-r68}
\end{equation}

\begin{proposition}[费米 Gaussian 态的四点配对]
设奇场满足 CAR，态为零均值准自由态；乘积顺序固定，配对符号由奇置换决定。则有以下结论（原式 \eqnrefs{eq:fermion-wick-four-r68}）：
\begin{equation*}
 \boxed{
 \langle0|\Tord\{\psi_1\bar\psi_2\psi_3\bar\psi_4\}|0\rangle
 =S_{12}S_{34}-S_{14}S_{32},}
\end{equation*}
\end{proposition}
\begin{proof}
正文\eqnrefs{eq:generic-fermion-contraction-dp8} 定义了自由费米场的一次收缩。对两对费米场只有直接配对和交叉配对两种完全收缩，而后者相对前者多一次 Grassmann 奇置换；把两项相加便得到正文\eqnrefs{eq:fermion-wick-four-r68}。



考虑固定初始顺序
\begin{equation*}
 \mathscr W_4
 =\langle0|T\{\psi_1\bar\psi_2\psi_3\bar\psi_4\}|0\rangle.
\end{equation*}
自由真空期望值必须把每个 $\psi$ 与一个 $\bar\psi$ 配对。第一种配对是 $(1,2)(3,4)$；两对在原顺序中已经相邻，因此不需要额外奇置换，贡献
\begin{equation*}
 +S_{12}S_{34}.
\end{equation*}
第二种配对是 $(1,4)(3,2)$。$\bar\psi_4$ 与 $\psi_1$ 收缩时必须越过中间的费米算符；等价地说，从配对 $(1,2)(3,4)$ 变到交叉配对改变了一次配对排列的宇称，因此贡献带负号。于是
\begin{equation*}
 {
 \mathscr W_4=S_{12}S_{34}-S_{14}S_{32}.}
 \end{equation*}
这个负号不是某个 QED 顶角的性质；即使没有电磁相互作用，自由费米高斯也已经具有它。

一般地，考虑 $n$ 个 $\psi_i$ 与 $n$ 个 $\bar\psi_j$。每个完全收缩等价于一个置换 $\pi\in S_n$，把第 $i$ 个 $\psi_i$ 配到 $\bar\psi_{\pi(i)}$。每次交换 Grassmann-odd 算符都产生 $-1$，故
\begin{equation*}
 \langle0|T\{\psi_1\bar\psi_1\cdots\psi_n\bar\psi_n\}|0\rangle
 =\sum_{\pi\in S_n}
 \operatorname{sgn}(\pi)
 \prod_{i=1}^{n}S_{i,\pi(i)}.
 \end{equation*}
右边正是矩阵 $S_{ij}$ 的行列式：
\begin{equation*}
 {
 \langle T\,\psi_1\bar\psi_1\cdots\psi_n\bar\psi_n\rangle_0
 =\det[S_{ij}].}
 \end{equation*}
与此相对，玻色高斯的完全收缩是不含 $\operatorname{sgn}(\pi)$ 的积和式配对和。Grassmann 高斯积分产生行列式，普通高斯积分产生行列式的负幂；它们与这里的 Wick 组合计数是同一反对易结构的不同表现。
\end{proof}

其中 $S_{ij}\equiv\langle0|\Tord\{\psi_i\bar\psi_j\}|0\rangle$。

在微扰展开中，若一组费米缩并首尾相接形成内部闭环，把该闭环上的 Grassmann 场重排为一个有序矩阵迹时会比对应的玻色配对多一次奇置换，因此每个独立闭合费米子圈带有
\begin{equation}
 \boxed{\text{each closed fermion loop contributes an overall factor }-1.}
 \label{eq:generic-closed-fermion-loop-r68}
\end{equation}

\begin{proposition}[闭费米收缩环的符号]
在 CAR 和固定时序的 Wick 展开中，比较具有同一外线顺序的完整闭合费米收缩链。则有以下结论（原式 \eqnrefs{eq:generic-closed-fermion-loop-r68}）：
\begin{equation*}
 \boxed{\text{each closed fermion loop contributes an overall factor }-1.}
\end{equation*}
\end{proposition}
\begin{proof}
在相互作用展开中，若若干费米收缩连成闭环，例如取如下示意性的双线性乘积

\begin{equation*}
 (\bar\psi_1\Gamma_1\psi_1)
 (\bar\psi_2\Gamma_2\psi_2)\cdots
 (\bar\psi_n\Gamma_n\psi_n),
\end{equation*}
并取循环收缩
\begin{equation*}
 \psi_1\leftrightarrow\bar\psi_2,
 \quad
 \psi_2\leftrightarrow\bar\psi_3,
 \quad\ldots\quad
 \psi_n\leftrightarrow\bar\psi_1,
\end{equation*}
则配对置换是一个长度为 $n$ 的循环。长度为 $n$ 的循环可分解成 $n-1$ 次对换，所以其宇称为 $(-1)^{n-1}$。另一方面，把每个双线性型中的 $\bar\psi_i$ 移到用于写矩阵迹的统一位置还产生 $(-1)^n$。两者相乘为
\begin{equation*}
 (-1)^{n-1}(-1)^n=-1.
\end{equation*}
因此一个闭合费米子圈相对于同样连接结构的对易的场多出
\begin{equation*}
 {-1.}
 \end{equation*}
\chapref{chap:electromagnetic-scattering}的 QED 真空极化的迹
$-\operatorname{Tr}[\gamma^\mu S_F\gamma^\nu S_F]$
就是这个一般结论的具体实例。
\end{proof}

具体相互作用顶角只决定环中的矩阵、内部量子数与传播子；这个整体负号在进入 QED、QCD 等具体理论之前已经由费米反对易代数固定。
